\documentclass[11pt,a4paper]{article}

% =========================================================
% COMMON PACKAGES
% =========================================================
\usepackage{lmodern}
\usepackage[T1]{fontenc}
\usepackage[utf8]{inputenc}
\usepackage[english]{babel}

\usepackage{amsmath,amssymb,amsfonts,amsthm,mathtools}
\usepackage{bm}
\usepackage{microtype}
\usepackage{enumitem}
\usepackage{csquotes}
\usepackage{mathrsfs}
\usepackage{thmtools}
\usepackage{hyperref}
\usepackage[nameinlink,capitalize]{cleveref}
\usepackage{geometry}

\geometry{
	a4paper,
	margin=2.8cm
}

% =========================================================
% THEOREM STYLES
% =========================================================
\numberwithin{equation}{section}

\declaretheoremstyle[
headfont=\bfseries,
bodyfont=\itshape,
spaceabove=8pt,
spacebelow=8pt
]{plainstyle}

\declaretheoremstyle[
headfont=\bfseries,
bodyfont=\normalfont,
spaceabove=8pt,
spacebelow=8pt
]{defstyle}

\declaretheoremstyle[
headfont=\itshape,
bodyfont=\normalfont,
spaceabove=4pt,
spacebelow=4pt
]{remarkstyle}

\declaretheorem[style=plainstyle,name=Theorem]{theorem}
\declaretheorem[style=plainstyle,name=Proposition,sibling=theorem]{proposition}
\declaretheorem[style=plainstyle,name=Lemma,sibling=theorem]{lemma}
\declaretheorem[style=plainstyle,name=Corollary,sibling=theorem]{corollary}
\declaretheorem[style=plainstyle,name=Conjecture,sibling=theorem]{conjecture}

\declaretheorem[style=defstyle,name=Definition,sibling=theorem]{definition}
\declaretheorem[style=defstyle,name=Example,sibling=theorem]{example}

\declaretheorem[style=remarkstyle,name=Remark,sibling=theorem]{remark}

% Cleveref names
\crefname{theorem}{Theorem}{Theorems}
\crefname{proposition}{Proposition}{Propositions}
\crefname{lemma}{Lemma}{Lemmas}
\crefname{corollary}{Corollary}{Corollaries}
\crefname{conjecture}{Conjecture}{Conjectures}
\crefname{definition}{Definition}{Definitions}
\crefname{example}{Example}{Examples}
\crefname{remark}{Remark}{Remarks}
\crefname{equation}{Equation}{Equations}
\crefname{section}{Section}{Sections}

% =========================================================
% COMMON MACROS FOR THE TWIN-SPACES SERIES
% =========================================================

% Sets and fields
\newcommand{\R}{\mathbb{R}}
\newcommand{\C}{\mathbb{C}}
\newcommand{\Q}{\mathbb{Q}}
\newcommand{\Z}{\mathbb{Z}}
\newcommand{\N}{\mathbb{N}}

% Operators
\DeclareMathOperator{\Spec}{Spec}
\DeclareMathOperator{\Proj}{Proj}
\DeclareMathOperator{\Pic}{Pic}
\DeclareMathOperator{\Divv}{Div}
\DeclareMathOperator{\Cl}{Cl}
\DeclareMathOperator{\CH}{CH}
\DeclareMathOperator{\Hom}{Hom}
\DeclareMathOperator{\Ker}{Ker}
\DeclareMathOperator{\coker}{coker}
\DeclareMathOperator{\Img}{Im}
\DeclareMathOperator{\Ext}{Ext}
\DeclareMathOperator{\Tor}{Tor}
\DeclareMathOperator{\End}{End}
\DeclareMathOperator{\Gal}{Gal}
\DeclareMathOperator{\Aut}{Aut}
\DeclareMathOperator{\id}{id}
\DeclareMathOperator{\Tr}{Tr}
\DeclareMathOperator{\degd}{deg}
\DeclareMathOperator{\codim}{codim}
\DeclareMathOperator{\Supp}{Supp}
\DeclareMathOperator{\rk}{rk}

% Twin notation
\newcommand{\TwinRing}{(R,G)}
\newcommand{\TwinScheme}{(X,G)}
\newcommand{\Twin}{\mathcal{T}}
\newcommand{\Ocal}{\mathcal{O}}
\newcommand{\Fcal}{\mathcal{F}}
\newcommand{\Gcal}{\mathcal{G}}
\newcommand{\Lcal}{\mathcal{L}}
\newcommand{\Ecal}{\mathcal{E}}

% Misc
\newcommand{\abs}[1]{\left\lvert #1 \right\rvert}
\newcommand{\norm}[1]{\left\lVert #1 \right\rVert}
\newcommand{\bracket}[1]{\left\langle #1 \right\rangle}

% =========================================================
% METADATA FOR ARTICLE D
% =========================================================

\title{\textbf{Twin Commutative Algebra and Algebraic Geometry:\\
		Schemes with Dynamic Group Actions and the Geometry of Symmetry}}

\author{Jocelyn Nembe}
\date{\today}

% =========================================================
% DOCUMENT
% =========================================================

\begin{document}
	
	\maketitle
	
	\begin{abstract}
		This article develops the algebraic–geometric layer of the twin spaces program.
		Building on the deformation cohomology theory (Article~A), the 2-categorical and operadic framework
		(Article~B), and the rigidity/quantum aspects (Article~C), we construct a full-fledged theory of 
		\emph{twin commutative algebra and algebraic geometry}.
		
		A twin ring is a pair $(R,G)$ where $R$ is a commutative ring and $G$ acts on $R$ by automorphisms.
		From this we derive twin ideals, the twin spectrum, and twin modules.
		A twin scheme is a scheme equipped with a compatible group action (or group scheme over the base), 
		and twin sheaves are sheaves endowed with equivariant structures.
		We introduce twin cohomology functors, a twin intersection theory, and a Riemann--Roch theorem for 
		twin varieties. We then construct twin moduli spaces (twin curves, twin Hilbert schemes), and develop 
		applications to public-key cryptography, error-correcting codes, and arithmetic dynamics.
		
		On the arithmetic side, we define twin \'etale cohomology and formulate a twin version of the Langlands program
		in which Galois representations and automorphic objects carry an additional twin structure.
		We conclude with examples (twin Fermat varieties, twin K3 surfaces), computational aspects
		(twin Gr\"obner bases and twin cohomology algorithms), and a list of open problems, including connections 
		with physics and motivic integration.
		
		Conceptually, twin algebraic geometry extends Grothendieck’s viewpoint by making symmetries 
		part of the foundational data of schemes, rather than an auxiliary structure added afterward.
	\end{abstract}
	
	\tableofcontents
	
	%===========================================================
	\section{Introduction}
	\label{sec:intro}
	
	\subsection{Classical background}
	
	Classical algebraic geometry studies solutions of polynomial equations via commutative algebra.
	Affine varieties correspond to prime spectra of rings, projective varieties to homogeneous spectra,
	and more generally schemes allow arithmetic and geometric information to be handled simultaneously.
	
	The Grothendieck paradigm may be summarized as:
	\begin{itemize}
		\item geometry is encoded by pairs $(X,\Ocal_X)$,
		\item sheaves and their cohomology carry refined information,
		\item moduli spaces organize families of geometric objects.
	\end{itemize}
	
	In many natural situations, however, the objects of interest are not just schemes, but schemes with 
	\emph{symmetry}:
	\begin{itemize}
		\item algebraic curves over finite fields come with Frobenius actions,
		\item complex varieties with real structures carry complex conjugation,
		\item abelian varieties have rich endomorphism rings and automorphism groups,
		\item moduli problems often have non-trivial stabilizers and monodromy.
	\end{itemize}
	
	Equivariant algebraic geometry and the theory of group actions on varieties address some of these phenomena,
	but the symmetry typically appears as an \emph{extra layer} on top of a given scheme.
	
	\subsection{The twin extension}
	
	The philosophy of twin spaces is that symmetry should be present \emph{from the beginning}:
	the underlying object and its symmetries form an inseparable pair.
	
	\begin{definition}[Twin ring, first informal version]
		A twin ring is a pair $(R,G)$ where $R$ is a commutative ring with identity and $G$ is a group
		acting on $R$ by ring automorphisms.
	\end{definition}
	
	This leads naturally to twin spectra, twin schemes, and twin sheaves.
	
	\begin{remark}[Twin perspective]
		A scheme encodes a topological space together with a sheaf of rings.
		A twin scheme enhances this by remembering also a group action (possibly varying over the base)
		and how this action interacts with the local algebra.
		This is particularly relevant for:
		\begin{enumerate}[label=(\alph*)]
			\item objects with rich intrinsic symmetry (e.g.\ modular curves, Shimura varieties),
			\item families with monodromy or varying symmetries,
			\item situations where symmetry-breaking and deformations are essential (physics, dynamics).
		\end{enumerate}
	\end{remark}
	
	\begin{example}[Natural sources of twin structure]
		Typical examples include:
		\begin{itemize}
			\item algebraic curves $C/\mathbb{F}_q$ with Frobenius endomorphism and $\Gal(\overline{\mathbb{F}}_q/\mathbb{F}_q)$,
			\item varieties over $\R$ with complex conjugation action of $\Z/2\Z$,
			\item abelian varieties with an action of an endomorphism ring or of a finite group,
			\item Calabi--Yau varieties with symmetry groups relevant in string theory.
		\end{itemize}
	\end{example}
	
	\subsection{Position in the twin spaces program}
	
	This article assumes the deformation and cohomological framework of twin operations developed in 
	Article~A, the 2-categorical and operadic structure of twin spaces in Article~B, and the rigidity/quantum
	aspects of Article~C. The present work is the \emph{algebraic geometry layer} of the program.
	
	Two features distinguish this layer from classical equivariant geometry. First, the invariant subring $R^G$
	and the invariant cohomology $H^i(X,\Fcal)^G$ are not afterthoughts but the primary global invariants of a
	twin object; when $\abs{G}$ is invertible they are exact functors (\cref{lem:averaging}) that behave like
	ordinary cohomology, whereas in the modular case they must be derived. Second, numerical invariants such as
	the Euler characteristic become \emph{averages over the group}, as the twin Riemann--Roch theorem
	(\cref{thm:RR-twin}) makes precise. The slogan throughout is that a twin scheme is a scheme whose symmetry is
	part of its definition, so every classical construction is replaced by its $G$-equivariant, then
	$G$-invariant, counterpart.

	The plan is as follows.
	In \cref{sec:twin-rings} we define twin rings and modules.
	In \cref{sec:twin-schemes} we introduce twin schemes and the twin spectrum functor.
	In \cref{sec:twin-sheaves} we develop twin sheaves and twin (co)homology.
	In \cref{sec:intersection} we set up twin intersection theory and twin Chow groups.
	In \cref{sec:riemann-roch} we prove a twin Riemann--Roch theorem.
	In \cref{sec:moduli} we construct twin moduli spaces.
	In \cref{sec:crypto} we discuss cryptographic applications.
	In \cref{sec:etale} we define twin \'etale cohomology.
	In \cref{sec:langlands} we sketch a twin Langlands picture.
	In \cref{sec:computational} we address computational issues.
	In \cref{sec:examples} we present examples.
	We conclude in \cref{sec:future} with open problems and connections to physics.
	
	%===========================================================
	\section{Twin Rings and Twin Modules}
	\label{sec:twin-rings}
	
	\subsection{Basic definitions}
	
	\begin{definition}[Twin ring]
		A \emph{twin ring} is a pair $(R,G)$ where:
		\begin{itemize}
			\item $R$ is a commutative ring with $1$,
			\item $G$ is a group,
			\item $G$ acts on $R$ by ring automorphisms
			\[
			\rho:G\times R\to R,\qquad (g,r)\mapsto g\cdot r,
			\]
			such that
			\[
			g\cdot(r+s)=g\cdot r+g\cdot s,\quad
			g\cdot(rs)=(g\cdot r)(g\cdot s),\quad
			g\cdot 1=1.
			\]
		\end{itemize}
	\end{definition}
	
	We write $R^G$ for the invariant subring, i.e.\ $R^G=\{r\in R : g\cdot r=r,\ \forall g\in G\}$.
	
	\begin{example}[Symmetric polynomials]
		Let $R=k[x_1,\dots,x_n]$ and $G=S_n$ acting by permutation of variables.
		Then $(R,S_n)$ is a twin ring and $R^{S_n}$ is the ring of symmetric polynomials. By the fundamental theorem of
		symmetric polynomials, $R^{S_n}=k[e_1,\dots,e_n]$ is itself a polynomial ring in the elementary symmetric
		functions, so $\Spec R^{S_n}\cong\mathbb A^n$ and the quotient map
		$\mathbb A^n=\Spec R\to\Spec R^{S_n}=\mathbb A^n$, $(x_1,\dots,x_n)\mapsto(e_1,\dots,e_n)$, is finite flat of
		degree $n!$, branched over the discriminant locus. This is the prototype twin spectrum: the geometry of the
		quotient encodes precisely the $S_n$-symmetry.
	\end{example}
	
	\begin{example}[Cyclotomic extension]
		Let $R=\Z[\zeta_n]$ where $\zeta_n$ is a primitive $n$-th root of unity, and let 
		$G=(\Z/n\Z)^\times$ act by $a\cdot\zeta_n=\zeta_n^a$.
		Then $(R,G)$ is a twin ring; $R^G=\Z$ reflects the Galois structure of the cyclotomic field.
	\end{example}
	
	\subsection{Twin modules}
	
	\begin{definition}[Twin module]
		Let $(R,G)$ be a twin ring.
		A \emph{twin $R$-module} is an $R$-module $M$ together with an action of $G$ by 
		$R$-linear automorphisms
		\[
		G\times M\to M,\quad (g,m)\mapsto g\cdot m
		\]
		such that
		\[
		g\cdot(rm) = (g\cdot r)(g\cdot m),
		\]
		for all $g\in G$, $r\in R$, $m\in M$.
	\end{definition}
	
	\begin{remark}
		Twin modules are precisely $R\rtimes G$-modules under the usual skew group ring construction.
		This is the algebraic shadow of the twin operations viewpoint of Articles~A~and~B.
	\end{remark}
	
	\subsection{The invariant subring and twin ideals}
	
	\begin{definition}[Invariant subring and action kernel]
		Let $(R,G)$ be a twin ring.
		The \emph{invariant subring} is
		\[
		R^G := \{\, r\in R : g\cdot r = r\ \text{for all } g\in G\,\}.
		\]
		It is a subring of $R$ containing $1$; in general it is \emph{not} an ideal of $R$ (an ideal
		containing $1$ would be all of $R$). The \emph{kernel} of the action is the normal subgroup
		\[
		N := \{\, g\in G : g\cdot r = r\ \text{for all } r\in R\,\} \trianglelefteq G.
		\]
	\end{definition}
	
	\begin{proposition}
		$R^G$ is a subring of $R$, and the action of $G$ on $R$ factors through a \emph{faithful} action
		of the quotient group $G/N$; in particular $N \neq G$ whenever the action is non-trivial. The
		inclusion $R^G \hookrightarrow R$ induces the quotient morphism $\Spec R \to \Spec R^G$.
	\end{proposition}
	
	\begin{proof}
		$R^G$ is closed under addition and multiplication and contains $1$, hence is a subring. (It is not
		an ideal: for $r\in R^G$ and $s\in R$ one has $g\cdot(sr)=(g\cdot s)r$, which differs from $sr$ as
		soon as $g\cdot s\neq s$, so $sr\notin R^G$ in general.) The assignment $g\mapsto(r\mapsto g\cdot r)$
		is a homomorphism $G\to\Aut(R)$ with kernel exactly $N$; hence $G/N$ acts faithfully, and $N\neq G$
		precisely when the action is non-trivial.
	\end{proof}
	
	\begin{definition}[Twin ideal]
		An ideal $I\subset R$ is called a \emph{twin ideal} if it is stable under $G$:
		\[
		g\cdot I \subseteq I\quad\text{for all } g\in G.
		\]
		If equality holds for all $g$, we call $I$ \emph{$G$-invariant}.
	\end{definition}
	
	\begin{remark}
		Twin ideals are the natural objects to consider for defining sub-twin rings, quotients, and spectra with symmetry.
		They form a lattice stable under finite intersections and sums.
	\end{remark}
	
	\subsection{Prime twin ideals and orbit structure}
	
	\begin{definition}[Twin prime]
		A proper \emph{prime} ideal $P\subset R$ is called \emph{twin-prime} if:
		\begin{itemize}
			\item $P$ is a prime ideal of $R$,
			\item its $G$-orbit $\{g\cdot P : g\in G\}$ is finite.
		\end{itemize}
	\end{definition}
	
	\begin{remark}
		The finiteness of the orbit ensures that the corresponding points in the spectrum form 
		a controlled configuration under the group action.
		For infinite groups one often restricts attention to orbits of finite type or closed orbits.
	\end{remark}
	
	\begin{proposition}
		Let $(R,G)$ be Noetherian.
		Then:
		\begin{enumerate}[label=(\alph*)]
			\item the set of twin-prime ideals is stable under the $G$-action,
			\item every twin ideal can be written as an intersection of twin-prime ideals (possibly after saturation).
		\end{enumerate}
	\end{proposition}
	
	\begin{proof}[Idea of proof]
		For (a), stability under $G$ is immediate. For (b), use the classical primary decomposition for $G$-invariant ideals 
		and group the components along $G$-orbits, taking intersections of the conjugates.
		Details follow standard arguments in equivariant algebra with the twin constraints built in.
	\end{proof}
	
	%===========================================================
	\section{Twin Schemes}
	\label{sec:twin-schemes}
	
	\subsection{Affine twin schemes}
	
	\begin{definition}[Affine twin scheme]
		Let $(R,G)$ be a twin ring.
		The \emph{affine twin scheme} associated to $\TwinRing$ is the pair
		\[
		\TwinScheme := (\Spec R, G),
		\]
		where $G$ acts on $\Spec R$ by
		\[
		g\cdot\mathfrak{p} := \{\, g\cdot r : r\in\mathfrak{p}\,\}
		\]
		and on the structure sheaf $\Ocal_{\Spec R}$ via the induced action on localizations.
	\end{definition}
	
	\begin{remark}
		The action on the structure sheaf is compatible with restriction and localization, so we obtain
		an action by automorphisms of ringed spaces. Affine twin schemes are thus objects of the 
		2-category of ringed spaces with group actions.
	\end{remark}
	
	\subsection{Twin spectrum functor}
	
	\begin{definition}[Twin spectrum]
		The \emph{twin spectrum} functor
		\[
		\Spec^{\mathrm{tw}} : (\text{Twin rings})^{\mathrm{op}} \to \text{Twin schemes}
		\]
		sends $(R,G)$ to $(\Spec R,G)$ and a morphism of twin rings
		\[
		\phi:(R,G)\to (S,H)
		\]
		(consisting of a ring map $\phi_R:R\to S$ and a group homomorphism $\phi_G:G\to H$
		compatible with the actions) to the twin morphism
		\[
		(\Spec \phi_R,\ \phi_G^\ast) : (\Spec S,H)\to(\Spec R,G),
		\]
		where $\Spec\phi_R$ is the usual map of schemes and $\phi_G^\ast$ is the induced action on group data.
	\end{definition}
	
	\begin{proposition}
		$\Spec^{\mathrm{tw}}$ is functorial, and restricts to the usual spectrum functor on the underlying rings.
	\end{proposition}
	
	\begin{proof}
		Functoriality is checked componentwise on rings and groups. Compatibility of the group action with the
		inverse image of primes ensures that the twin structure is preserved.
	\end{proof}
	
	\subsection{Gluing twin affine pieces}
	
	\begin{definition}[Twin scheme]
		A \emph{twin scheme} is a pair $(X,G)$ where:
		\begin{itemize}
			\item $X$ is a scheme,
			\item $G$ is a group (or group scheme) acting on $X$ by automorphisms,
			\item there exists an open affine cover $\{U_i\}$ of $X$ such that each $(U_i,G)$ is isomorphic
			to an affine twin scheme $(\Spec R_i,G)$ for some twin ring $(R_i,G)$, and the gluing maps
			are $G$-equivariant.
		\end{itemize}
	\end{definition}
	
	\begin{proposition}
		Twin schemes and twin morphisms (scheme morphisms compatible with the group actions) form a category,
		and affine twin schemes are precisely the objects isomorphic to $\Spec^{\mathrm{tw}}(R,G)$.
	\end{proposition}
	
	\begin{proof}
		Gluing along $G$-equivariant isomorphisms is standard in descent theory for schemes.
		The compatibility conditions ensure that the resulting action is globally defined and the structure sheaf 
		is preserved.
	\end{proof}
	
	\begin{example}[Twin projective space]
		Let $R=k[x_0,\dots,x_n]$ with $G\subset\Aut_k(R)$ a group of homogeneous automorphisms.
		Then $X=\Proj R$ carries an induced $G$-action, and $(\Proj R,G)$ is a twin scheme.
	\end{example}
	
	%===========================================================
	\section{Twin Sheaves and Twin Cohomology}
	\label{sec:twin-sheaves}
	
	\subsection{Twin sheaves}
	
	\begin{definition}[Twin $\Ocal_X$-module]
		Let $\TwinScheme$ be a twin scheme.
		A \emph{twin sheaf} (or twin $\Ocal_X$-module) on $\TwinScheme$ is a sheaf of $\Ocal_X$-modules $\Fcal$
		together with, for each $g\in G$, an isomorphism of sheaves
		\[
		\theta_g : g^\ast\Fcal \xrightarrow{\sim} \Fcal
		\]
		such that $\theta_{1}=\id$ and $\theta_{gh}=\theta_h\circ h^\ast\theta_g$ for all $g,h\in G$.
	\end{definition}
	
	\begin{remark}
		This is the usual notion of a $G$-equivariant sheaf, but in the twin viewpoint it is part of the
		intrinsic structure: twin schemes are always considered together with their twin sheaves.
	\end{remark}
	
	\begin{definition}[Twin quasi-coherent and coherent sheaves]
		A twin sheaf $\Fcal$ is called \emph{twin quasi-coherent} (resp.\ \emph{twin coherent}) if the 
		underlying $\Ocal_X$-module is quasi-coherent (resp.\ coherent).
	\end{definition}
	
	\subsection{Twin cohomology}
	
	\begin{definition}[Twin cohomology]
		Let $\Fcal$ be a twin sheaf on $\TwinScheme$.
		The \emph{twin cohomology} of $\Fcal$ is defined as
		\[
		H^i_{\mathrm{tw}}(X,\Fcal) := H^i(X,\Fcal)^G,
		\]
		the $G$-invariant part of the usual sheaf cohomology.
	\end{definition}
	
	\begin{remark}
		One can refine this by considering the full representation $H^i(X,\Fcal)$ as a $G$-module and
		analyzing its decomposition, but for the present article we focus on the invariant part, which is the
		natural receptacle for global twin structures.
	\end{remark}
	
	\begin{lemma}[Averaging and exactness of invariants]\label{lem:averaging}
		Suppose $G$ is finite and $\abs{G}$ is invertible in the base field (e.g.\ characteristic zero). Then the
		Reynolds operator $e=\tfrac1{\abs G}\sum_{g\in G}g$ is an idempotent endomorphism of every $G$-module $M$
		with image $M^G$, giving a canonical splitting $M=M^G\oplus(1-e)M$. Hence the functor $(-)^G$ is exact, and
		$H^i_{\mathrm{tw}}(X,\Fcal)=H^i(X,\Fcal)^G$ is a direct summand of $H^i(X,\Fcal)$ on which twin cohomology
		inherits the long exact sequence of any short exact sequence of twin sheaves. In the \emph{modular} case
		($\operatorname{char}\mid\abs G$) exactness fails and the correct object is the $G$-hypercohomology,
		computed by a spectral sequence $H^p\!\big(G,H^q(X,\Fcal)\big)\Rightarrow \mathbb H^{p+q}$.
	\end{lemma}

	\begin{proof}
		Idempotency $e^2=e$ and $eM=M^G$ are the standard properties of the Reynolds operator; $1-e$ projects onto a
		$G$-stable complement. A functor that is a direct-summand projection is exact, and exact functors commute
		with the formation of cohomology of complexes, giving the stated long exact sequence. The modular statement
		is the Hochschild--Serre/equivariant hypercohomology spectral sequence.
	\end{proof}

	\begin{proposition}[Functoriality]
		Twin cohomology is functorial in the pair $(X,\Fcal)$ and compatible with pullback and pushforward
		along twin morphisms under standard hypotheses (properness, flatness, etc.).
	\end{proposition}
	
	\begin{proof}[Idea]
		Functoriality of usual cohomology together with the equivariance of morphisms of twin schemes and twin sheaves
		implies that the induced maps preserve $G$-invariants.
	\end{proof}
	
	\subsection{Twin \v{C}ech cohomology}
	
	Let $\mathfrak{U}=\{U_i\}$ be a $G$-stable open cover of $X$, i.e.\ $gU_i$ is a union of $U_j$'s for all $g\in G$.
	
	\begin{definition}[Twin \v{C}ech complex]
		The \emph{twin \v{C}ech complex} associated with $\mathfrak{U}$ and a twin sheaf $\Fcal$ is the usual
		\v{C}ech complex
		\[
		\check{C}^\bullet(\mathfrak{U},\Fcal),
		\]
		viewed as a complex of $G$-modules via the twin structure. Its $G$-invariant subcomplex
		\[
		\check{C}^\bullet_{\mathrm{tw}}(\mathfrak{U},\Fcal) := \check{C}^\bullet(\mathfrak{U},\Fcal)^G
		\]
		computes the twin \v{C}ech cohomology
		\[
		\check{H}^i_{\mathrm{tw}}(\mathfrak{U},\Fcal) := H^i\big(\check{C}^\bullet_{\mathrm{tw}}(\mathfrak{U},\Fcal)\big).
		\]
	\end{definition}
	
	\begin{theorem}[Comparison]
		If $\Fcal$ is twin quasi-coherent on a separated, quasi-compact twin scheme and $\mathfrak{U}$ is a 
		sufficiently fine $G$-stable affine cover, then
		\[
		\check{H}^i_{\mathrm{tw}}(\mathfrak{U},\Fcal) \cong H^i_{\mathrm{tw}}(X,\Fcal).
		\]
	\end{theorem}
	
	\begin{proof}[Sketch]
		A $G$-stable affine cover computes $H^i(X,\Fcal)$ because $X$ is separated and quasi-compact and $\Fcal$ is
		quasi-coherent. By the averaging \cref{lem:averaging}, when $\abs G$ is invertible the exact functor $(-)^G$
		commutes with the cohomology of the \v{C}ech complex, so
		$\check H^i_{\mathrm{tw}}(\mathfrak U,\Fcal)=\check H^i(\mathfrak U,\Fcal)^G\cong H^i(X,\Fcal)^G=H^i_{\mathrm{tw}}(X,\Fcal)$.
		In the modular case the same holds after replacing invariants by $G$-hypercohomology.
	\end{proof}
	
	%===========================================================
	\section{Twin Intersection Theory}
	\label{sec:intersection}
	
	\subsection{Twin divisors}
	
	Let $\TwinScheme$ be an integral twin scheme, normal and separated of finite type.
	
	\begin{definition}[Twin Weil divisor]
		A \emph{twin Weil divisor} on $\TwinScheme$ is a finite formal $\Z$-linear combination of 
		irreducible codimension-one subvarieties of $X$ that are $G$-stable, i.e.
		\[
		D = \sum_i n_i Z_i,\qquad n_i\in\Z,
		\]
		with each $Z_i$ invariant under the action of $G$.
		We write $\Divv_{\mathrm{tw}}(X)$ for the group of twin Weil divisors.
	\end{definition}
	
	\begin{definition}[Twin Cartier divisor]
		A \emph{twin Cartier divisor} is a Cartier divisor $D$ on $X$ represented by local equations
		compatible with the $G$-action, i.e.\ for each $g\in G$ the pullback $g^\ast D$ is linearly equivalent to $D$.
		We denote by $\Pic_{\mathrm{tw}}(X)$ the group of isomorphism classes of twin line bundles.
	\end{definition}
	
	\begin{remark}
		The twin Picard group is a subgroup of the usual Picard group $\Pic(X)$, 
		consisting of line bundles endowed with a $G$-equivariant structure.
	\end{remark}
	
	\subsection{Twin Chow groups}
	
	\begin{definition}[Twin cycles]
		For $k\ge 0$, a \emph{twin $k$-cycle} on $\TwinScheme$ is a formal linear combination (with integer coefficients)
		of irreducible closed subschemes of dimension $k$ that are stabilized by $G$.
		The group of twin $k$-cycles is denoted $Z_k^{\mathrm{tw}}(X)$.
	\end{definition}
	
	\begin{definition}[Twin rational equivalence]
		Two twin cycles are said to be twin-rationally equivalent if their difference is a finite sum of cycles of 
		the form $\mathrm{div}(f)$ where $f$ is a rational function compatible with the $G$-action.
		The quotient group
		\[
		\CH_k^{\mathrm{tw}}(X) := Z_k^{\mathrm{tw}}(X)/\sim_{\mathrm{tw}}
		\]
		is called the \emph{twin Chow group} of dimension $k$.
	\end{definition}
	
	\begin{remark}
		Equivalently, one can define $\CH_k^{\mathrm{tw}}(X)$ as the $G$-invariant part of the usual Chow group,
		endowed with the induced intersection product.
	\end{remark}
	
	\subsection{Intersection pairings}
	
	Under usual smoothness and properness assumptions, $\CH_*^{\mathrm{tw}}(X)$ carries an intersection product
	\[
	\cdot : \CH_i^{\mathrm{tw}}(X)\times \CH_j^{\mathrm{tw}}(X) \to \CH_{i+j-d}^{\mathrm{tw}}(X),
	\]
	where $d=\dim X$.
	
	\begin{proposition}
		The intersection product of cycles descends to twin Chow groups and is compatible with the $G$-action.
		In particular, the $G$-invariant part of the classical intersection pairing induces a well-defined pairing on
		$\CH_*^{\mathrm{tw}}(X)$.
	\end{proposition}
	
	\begin{proof}[Idea]
		Since the intersection product is functorial and $G$ acts by automorphisms of the scheme, 
		$g\cdot(Z\cdot Z')=(g\cdot Z)\cdot(g\cdot Z')$ for all cycles.
		Passing to $G$-invariant cycles and rational equivalence classes yields the result.
	\end{proof}
	
	%===========================================================
	\section{Twin Riemann--Roch Theorem}
	\label{sec:riemann-roch}
	
	We now outline a Riemann--Roch formula in the twin setting.
	
	\subsection{Twin Chern and Todd classes}
	
	Let $(X,G)$ be a smooth projective twin scheme over a field of characteristic zero, and let $\Ecal$ be a twin coherent
	sheaf (e.g.\ a twin vector bundle) on $X$.
	
	\begin{definition}[Twin Chern character]
		The \emph{twin Chern character} of $\Ecal$ is defined as
		\[
		\mathrm{ch}_{\mathrm{tw}}(\Ecal) := \mathrm{ch}(\Ecal)^G \in \bigoplus_i \CH^i_{\mathrm{tw}}(X)\otimes\Q,
		\]
		the $G$-invariant part of the usual Chern character.
	\end{definition}
	
	\begin{definition}[Twin Todd class]
		The \emph{twin Todd class} of $X$ is
		\[
		\mathrm{Td}_{\mathrm{tw}}(X) := \mathrm{Td}(T_X)^G \in \bigoplus_i \CH^i_{\mathrm{tw}}(X)\otimes\Q.
		\]
	\end{definition}
	
	\subsection{The twin Riemann--Roch statement}
	
	\begin{definition}[Twin Euler characteristic]
		The \emph{twin Euler characteristic} of $\Ecal$ is
		\[
		\chi_{\mathrm{tw}}(X,\Ecal) := \sum_{i\ge 0} (-1)^i \dim H^i_{\mathrm{tw}}(X,\Ecal).
		\]
	\end{definition}
	
	\begin{theorem}[Twin Grothendieck--Riemann--Roch]\label{thm:RR-twin}
		Let $(X,G)$ be a smooth projective twin scheme and $\Ecal$ a twin coherent sheaf.
		Then
		\[
		\chi_{\mathrm{tw}}(X,\Ecal) 
		= \frac{1}{\abs{G}}\sum_{g\in G}\chi(g,\Ecal),
		\qquad
		\chi(g,\Ecal):=\sum_{i\ge0}(-1)^i\operatorname{tr}\!\big(g\mid H^i(X,\Ecal)\big),
		\]
		where $\chi(g,\Ecal)$ is the $g$-equivariant Euler characteristic, computed by the holomorphic Lefschetz
		fixed-point formula of Atiyah--Bott as an integral over the fixed locus $X^g$.
	\end{theorem}
	
	\begin{proof}[Idea of proof]
		By the averaging \cref{lem:averaging}, for finite $G$ in characteristic zero
		$\dim H^i(X,\Ecal)^G=\tfrac1{\abs G}\sum_{g}\operatorname{tr}(g\mid H^i(X,\Ecal))$, the rank of the Reynolds
		idempotent. Summing with signs,
		\[
			\chi_{\mathrm{tw}}(X,\Ecal)=\sum_i(-1)^i\dim H^i(X,\Ecal)^G
			=\frac1{\abs G}\sum_{g\in G}\sum_i(-1)^i\operatorname{tr}(g\mid H^i)
			=\frac1{\abs G}\sum_{g\in G}\chi(g,\Ecal).
		\]
		Each $\chi(g,\Ecal)$ is given by the holomorphic Lefschetz (Atiyah--Bott) fixed-point theorem as an integral
		over $X^g$; the term $g=e$ contributes $\chi(e,\Ecal)=\chi(X,\Ecal)=\int_X\operatorname{ch}(\Ecal)\operatorname{Td}(X)$
		by ordinary Grothendieck--Riemann--Roch. This proves the averaged formula.
		
		 
		
		
	\end{proof}
	
	\begin{remark}
		In the presence of non-trivial stabilizers, local contributions from fixed-point loci of the $G$-action appear,
		leading to refined formulas reminiscent of the Atiyah--Segal--Singer index theorem in the equivariant setting.
	\end{remark}

	\begin{remark}[The naive product formula is \emph{not} twin Riemann--Roch]\label{rem:RR-warning}
		One might guess $\chi_{\mathrm{tw}}(X,\Ecal)=\int_X\operatorname{ch}_{\mathrm{tw}}(\Ecal)\cdot\operatorname{Td}_{\mathrm{tw}}(X)$
		with the invariant characteristic classes defined above. This is false in general, because the invariant
		part of a product is not the product of invariant parts. Take $X=\mathbb P^1$, $G=\Z/2$ acting by
		$[x_0:x_1]\mapsto[x_0:-x_1]$, and $\Ecal=\Ocal(1)$: then $H^0=\langle x_0,x_1\rangle$ has $\sigma$-eigenvalues
		$\pm1$, so $\dim(H^0)^G=1$ and $\chi_{\mathrm{tw}}=1$; but $G$ acts trivially on $\CH(\mathbb P^1)$, whence
		$\operatorname{ch}_{\mathrm{tw}}=1+[\mathrm{pt}]$, $\operatorname{Td}_{\mathrm{tw}}=1+[\mathrm{pt}]$, and
		$\int_{\mathbb P^1}(1+[\mathrm{pt}])^2=2\neq1$. The averaged formula gives the correct value
		$\tfrac12(\chi(e)+\chi(\sigma))=\tfrac12(2+0)=1$. The product formula holds only when $G$ acts trivially on
		every $H^i(X,\Ecal)$.
	\end{remark}
	
	%===========================================================
	\section{Twin Moduli Spaces}
	\label{sec:moduli}
	
	\subsection{Moduli of twin curves}
	
	Let $(C,G)$ be a smooth projective twin curve of genus $g\ge 2$ over an algebraically closed field $k$,
	where $G$ acts by automorphisms.
	
	\begin{definition}[Twin curve]
		A \emph{twin curve} of type $(g,G)$ is a pair $(C,G)$ where $C$ is a smooth projective curve of genus $g$
		and $G\subset\Aut(C)$ is a fixed finite group acting faithfully on $C$.
	\end{definition}
	
	\begin{definition}[Moduli functor]
		The moduli functor of twin curves assigns to a $k$-scheme $S$ the groupoid of families
		\[
		\pi:\mathcal{C}\to S
		\]
		of smooth proper curves of genus $g$ together with an action of the constant group scheme $G_S$ on $\mathcal{C}$
		over $S$, satisfying natural stability conditions.
	\end{definition}
	
	\begin{theorem}[Existence of moduli]
		Under standard stable reduction hypotheses, the moduli functor of twin curves of type $(g,G)$ is represented (or
		at least corepresented) by a Deligne--Mumford stack $\mathcal{M}_{g}^{\mathrm{tw}}(G)$ equipped with a natural 
		twin structure.
	\end{theorem}
	
	\begin{remark}
		The stack $\mathcal{M}_{g}^{\mathrm{tw}}(G)$ sits inside the usual moduli stack of curves as a closed substack
		parameterizing curves with specified automorphism group data. The twin viewpoint promotes this to the 
		\emph{primary} moduli object, not just a special locus.
	\end{remark}
	
	\subsection{Twin Hilbert schemes}
	
	\begin{definition}[Twin Hilbert functor]
		Given a twin scheme $(X,G)$ and a fixed Hilbert polynomial $P$, the \emph{twin Hilbert functor} sends a $k$-scheme $S$
		to the groupoid of closed subschemes $Z\subset X\times S$ flat over $S$ with Hilbert polynomial $P$ and stable under the
		action of $G$ (i.e.\ $G$ acts fiberwise on $Z$).
	\end{definition}
	
	\begin{theorem}[Twin Hilbert scheme]
		Under suitable Noetherian hypotheses, the twin Hilbert functor is represented by a twin scheme 
		$\mathrm{Hilb}^{\mathrm{tw}}_P(X,G)$, called the \emph{twin Hilbert scheme}.
	\end{theorem}
	
	\begin{proof}[Idea]
		One constructs the usual Hilbert scheme and then considers the closed subscheme where 
		the universal family is $G$-stable.
		The action of $G$ on $X$ induces an action on the Hilbert scheme, and the fixed locus 
		captures twin structures.
	\end{proof}
	
	%===========================================================
	\section{Applications to Cryptography and Coding Theory}
	\label{sec:crypto}
	
	\subsection{The hidden kernel problem}
	
	Twin rings and twin schemes give rise to new hardness assumptions for cryptography.
	
	\begin{definition}[Hidden kernel problem (HKP)]
		Let $(R,G)$ be a twin ring and $\varphi:R\to S$ a surjective ring homomorphism to another ring $S$
		carrying a compatible $G$-action. The \emph{hidden kernel problem} asks, given black-box access to $\varphi$
		and to the $G$-action, to recover the kernel ideal $\Ker\varphi$ (or its $G$-saturation).
	\end{definition}
	
	\begin{remark}
		The twin structure makes the kernel invariant under $G$ and often highly non-trivial combinatorially,
		suggesting potential post-quantum hardness when $G$ is large or when the geometry of $\Spec R$ is complicated.
	\end{remark}
	
	\subsection{Twin elliptic curves}
	
	\begin{example}[Twin elliptic curves]
		Let $E/k$ be an elliptic curve and $G$ a finite subgroup of $\Aut(E)$ (for example, complex multiplication
		or Frobenius over finite fields).
		Then $(E,G)$ is a twin curve, and one can define twin endomorphism rings, twin isogenies, and twin torsion points.
		These structures yield variants of isogeny-based cryptosystems where the extra symmetry may be 
		used either as a feature (extra structure for protocols) or as a parameter to hide.
	\end{example}
	
	\subsection{Twin algebraic codes}
	
	\begin{definition}[Twin algebraic--geometry code]
		Let $\TwinScheme$ be a twin curve over a finite field, $\Lcal$ a twin line bundle, and $P_1,\dots,P_n\in X(k)$
		a collection of rational points forming $G$-orbits. The associated \emph{twin algebraic--geometry code} is defined 
		by evaluating global twin sections
		\[
		H^0_{\mathrm{tw}}(X,\Lcal)\subset H^0(X,\Lcal)
		\]
		at the points $P_i$.
	\end{definition}
	
	\begin{remark}
		Twin codes inherit symmetry properties from $G$ that can be exploited to design codes with large automorphism groups,
		potentially useful for code-based cryptography and efficient decoding algorithms.
	\end{remark}
	
	%===========================================================
	\section{Twin \'Etale Cohomology}
	\label{sec:etale}
	
	\subsection{Twin \'etale site}
	
	Let $\TwinScheme$ be a twin scheme over a base where \'etale cohomology is well-behaved (e.g.\ Noetherian, finite type).
	
	\begin{definition}[Twin \'etale site]
		The \emph{twin \'etale site} $X_{\mathrm{\acute{e}t}}^{\mathrm{tw}}$ is the usual \'etale site of $X$ equipped with the 
		induced $G$-action on objects and morphisms.
		A \emph{twin \'etale sheaf} is an \'etale sheaf together with a compatible $G$-equivariant structure.
	\end{definition}
	
	\begin{definition}[Twin \'etale cohomology]
		For a twin \'etale sheaf $\Fcal$ on $X_{\mathrm{\acute{e}t}}^{\mathrm{tw}}$, the \emph{twin \'etale cohomology} is
		\[
		H^i_{\mathrm{\acute{e}t,tw}}(X,\Fcal) := H^i_{\mathrm{\acute{e}t}}(X,\Fcal)^G.
		\]
	\end{definition}
	
	\begin{remark}
		For constant coefficients $\Fcal=\underline{\Lambda}$ with a non-trivial $G$-action, this captures 
		simultaneously arithmetic information (via the \'etale site) and symmetry information.
	\end{remark}
	
	\subsection{$\ell$-adic twin cohomology}
	
	Let $\ell$ be a prime different from the residue characteristic.
	Take an inverse system of twin \'etale sheaves $(\Fcal_n)_{n\ge 1}$ with $\Fcal_n$ $\ell^n$-torsion and 
	transition maps compatible with $G$.
	
	\begin{definition}[$\ell$-adic twin cohomology]
		The \emph{$\ell$-adic twin cohomology} is defined as
		\[
		H^i_{\mathrm{\acute{e}t,tw}}(X,\Z_\ell) := \varprojlim_n H^i_{\mathrm{\acute{e}t}}(X,\Fcal_n)^G.
		\]
	\end{definition}
	
	\begin{remark}
		The resulting $\Z_\ell$-module carries an action of the absolute Galois group and of $G$, 
		and thus is a natural receptacle for twin Galois representations.
	\end{remark}
	
	%===========================================================
	\section{The Twin Langlands Program}
	\label{sec:langlands}
	
	\subsection{Twin Galois representations}
	
	Let $k$ be a global field and $X/k$ a smooth projective twin scheme.
	
	\begin{definition}[Twin Galois representation]
		An \emph{$\ell$-adic twin Galois representation} is a continuous representation
		\[
		\rho_{\mathrm{tw}} : \Gal(\overline{k}/k)\to \Aut_{\Z_\ell}\big(H^i_{\mathrm{\acute{e}t,tw}}(X_{\overline{k}},\Z_\ell)\big)
		\]
		compatible with the $G$-action.
	\end{definition}
	
	\begin{remark}
		Classical $\ell$-adic representations are recovered by forgetting the $G$-structure.
		The twin data encode additional symmetry coming from the twin scheme.
	\end{remark}
	
	\subsection{Twin automorphic forms}
	
	\begin{definition}[Twin automorphic object, very rough]
		A \emph{twin automorphic object} is an automorphic function, form, or representation on an adelic group
		equipped with a compatible action of a group $G$ of symmetries, arising from the twin geometry of 
		a twin scheme or a twin moduli space.
	\end{definition}
	
	\begin{conjecture}[Twin Langlands correspondence, schematic]
		There exists a correspondence between suitable classes of twin Galois representations and 
		twin automorphic objects, intertwining:
		\begin{itemize}
			\item the action of $\Gal(\overline{k}/k)$ on twin \'etale cohomology;
			\item the Hecke action and the twin symmetry $G$ on automorphic side.
		\end{itemize}
	\end{conjecture}
	
	\begin{remark}
		This is a long-term direction.
		Twin structures could provide a cleaner way to package internal symmetries of varieties 
		into the Langlands formalism, especially for moduli spaces with non-trivial automorphism groups.
	\end{remark}
	
	%===========================================================
	\section{Computational Aspects}
	\label{sec:computational}
	
	\subsection{Gr\"obner bases for twin ideals}
	
	\begin{definition}[Twin Gr\"obner basis]
		Let $(R,G)$ be a twin polynomial ring over a field, and $I\subset R$ a twin ideal.
		A \emph{twin Gr\"obner basis} of $I$ is a Gr\"obner basis $G_I$ that is stable under $G$:
		$g\cdot G_I=G_I$ for all $g\in G$.
	\end{definition}
	
	\begin{remark}
		Twin Gr\"obner bases reflect both algebraic and symmetry constraints and can drastically 
		reduce redundancy in computations, since one can work up to the $G$-action.
	\end{remark}
	
	\subsection{Computing twin cohomology}
	
	For twin sheaves on projective twin schemes, one can adapt standard algorithms for sheaf cohomology
	(complexes from free resolutions, \v{C}ech complexes) and then take $G$-invariants at the cochain level.
	
	\begin{proposition}[Algorithmic outline]
		Assuming effective presentations of $R$, $G$, and $\Fcal$, one can compute 
		$H^i_{\mathrm{tw}}(X,\Fcal)$ by:
		\begin{enumerate}[label=(\roman*)]
			\item constructing a cochain complex computing $H^i(X,\Fcal)$,
			\item describing the $G$-action on each degree,
			\item taking invariants degreewise and computing cohomology of the subcomplex.
		\end{enumerate}
	\end{proposition}
	
	\begin{remark}
		Practical implementations require careful control of complexity, but the presence of symmetry
		often reduces the size of the problem significantly.
	\end{remark}
	
	%===========================================================
	\section{Examples}
	\label{sec:examples}
	
	\subsection{Twin Fermat varieties}
	
	\begin{example}[Twin Fermat hypersurface]
		Let $X\subset\mathbb{P}^{n+1}$ be the Fermat hypersurface defined by
		\[
		x_0^{d}+\cdots+x_{n+1}^{d}=0
		\]
		over a field containing the $d$-th roots of unity, and let $G$ be the group of diagonal symmetries
		generated by multiplying coordinates by roots of unity with product $1$.
		Then $(X,G)$ is a twin variety with rich symmetry.
		Twin cohomology isolates the $G$-invariant part of the Hodge structure and is expected to relate
		to special values of $L$-functions and periods.
	\end{example}
	
	\subsection{Twin K3 surfaces}
	
	\begin{example}[Twin K3]
		Let $X$ be a K3 surface with a finite symplectic automorphism group $G$.
		Then $(X,G)$ is a twin K3 surface.
		Twin intersection theory on $\CH^1_{\mathrm{tw}}(X)$ captures the $G$-fixed part of the N\'eron--Severi group,
		and twin \'etale cohomology provides refined Galois representations with extra symmetry.
		Such objects are natural test beds for the twin Langlands conjectures.
	\end{example}
	
	%===========================================================
	\section{Future Directions and Connections to Physics}
	\label{sec:future}
	
	We conclude with a brief list of directions suggested by the preceding theory.
	
	\subsection{Open problems}
	
	\begin{itemize}
		\item \textbf{Twin resolution and minimal models.}
		Develop a theory of minimal models in the twin category, compatible with $G$-actions and twin cohomology.
		\item \textbf{Twin motives.}
		Construct a twin version of the category of motives where correspondences carry twin structure.
		\item \textbf{Twin Donaldson--Thomas invariants.}
		Define enumerative invariants for twin Calabi--Yau threefolds counting twin-stable objects.
		\item \textbf{Effective algorithms.}
		Provide complexity bounds for computing twin Hilbert schemes and twin cohomology for explicit examples.
	\end{itemize}
	
	\subsection{Connections to physics}
	
	Twin schemes appear naturally in contexts where symmetry plays a central r\^ole:
	\begin{itemize}
		\item \emph{String theory and mirror symmetry.}
		Calabi--Yau manifolds with large symmetry groups suggest studying twin versions of mirror symmetry.
		\item \emph{Gauge theory.}
		Moduli of bundles or sheaves with structure group $G$ over a twin base could reveal new dualities.
		\item \emph{Quantum field theory.}
		Twin operations and quantum twin spaces (Article~C) might be related to algebraic structures 
		underlying topological quantum field theories.
	\end{itemize}
	
	These perspectives reinforce the idea that twin algebraic geometry is not only a natural extension of 
	classical scheme theory, but also a bridge between algebraic geometry, representation theory, and 
	mathematical physics.
	
	%===========================================================
	\begin{thebibliography}{99}
		
		\bibitem{Hartshorne}
		R.~Hartshorne,
		\newblock {\em Algebraic Geometry},
		\newblock Graduate Texts in Mathematics, 52, Springer, 1977.
		
		\bibitem{EGA}
		A.~Grothendieck and J.~Dieudonn\'e,
		\newblock {\em \'El\'ements de G\'eom\'etrie Alg\'ebrique},
		\newblock Inst. Hautes \'Etudes Sci. Publ. Math., 1960--1967.
		
		\bibitem{SGA1}
		A.~Grothendieck,
		\newblock {\em Rev\^etements \'Etales et Groupe Fondamental (SGA 1)},
		\newblock Lecture Notes in Mathematics, 224, Springer, 1971.
		
		\bibitem{MumfordGIT}
		D.~Mumford, J.~Fogarty, and F.~Kirwan,
		\newblock {\em Geometric Invariant Theory},
		\newblock 3rd ed., Ergebnisse der Mathematik, Springer, 1994.
		
		\bibitem{Fulton}
		W.~Fulton,
		\newblock {\em Intersection Theory},
		\newblock 2nd ed., Ergebnisse der Mathematik, Springer, 1998.
		
		\bibitem{MilneEtale}
		J.~S.~Milne,
		\newblock {\em \'Etale Cohomology},
		\newblock Princeton Math. Series, 33, Princeton Univ. Press, 1980.
		
		\bibitem{SerreLocalFields}
		J.-P.~Serre,
		\newblock {\em Local Fields},
		\newblock Graduate Texts in Mathematics, 67, Springer, 1979.
		
		\bibitem{ChariPressley}
		V.~Chari and A.~Pressley,
		\newblock {\em A Guide to Quantum Groups},
		\newblock Cambridge Univ. Press, 1994.
		
		\bibitem{Weibel}
		C.~A.~Weibel,
		\newblock {\em An Introduction to Homological Algebra},
		\newblock Cambridge Studies in Advanced Mathematics, 38, Cambridge Univ. Press, 1994.
		
		\bibitem{NembeArticleI}
		J.~Nembe,
		\newblock Twin Operations: A Universal Framework for Algebraic Deformation via Group Actions,
		\newblock Article I in this series, 2025.
		
		\bibitem{NembeArticleII}
		J.~Nembe,
		\newblock Twin Operations II: Cohomological and Categorical Perspectives,
		\newblock Article II in this series, 2025.
		
		\bibitem{NembeArticleIII}
		J.~Nembe,
		\newblock Rigidity, Flexibility, and Quantum Deformations of Twin Spaces,
		\newblock Article III in this series, 2025.
		
	\end{thebibliography}
	
\end{document}
